\documentclass[10pt,a4paper]{amsart}
\usepackage{amsmath,amssymb,amsthm,graphicx}
% Retained minimal preamble for the Stacks Project chapter more-algebra.tex (More on Algebra).
% The chapter file pinned-original-source.tex carries only \input{preamble} on its first line, so no
% preamble text lives inside the chapter itself.  The Stacks Project uses one single shared preamble for
% every chapter; the snapshot stages it identically as sources/stacks-chow/preamble.tex and
% sources/stacks-etale-cohomology/preamble.tex.  The block below is copied mechanically, line for line,
% from sources/stacks-chow/preamble.tex lines 14-18 (the Xy-pic setup) and 154-217 (the native theorem
% declarations together with the mathematical macros).  Nothing was added, removed or redefined; the
% wrapper class and the amsmath/amssymb/amsthm/graphicx base packages are supplied by the assembler.
% Chapter: more-algebra.tex -- More on Algebra.
% For commutative diagrams we use Xy-pic
\usepackage[all]{xy}

% We use 2cell for 2-commutative diagrams.
\xyoption{2cell}

% Theorem environments.
%
\theoremstyle{plain}
\newtheorem{theorem}[subsection]{Theorem}
\newtheorem{proposition}[subsection]{Proposition}
\newtheorem{lemma}[subsection]{Lemma}

\theoremstyle{definition}
\newtheorem{definition}[subsection]{Definition}
\newtheorem{example}[subsection]{Example}
\newtheorem{exercise}[subsection]{Exercise}
\newtheorem{situation}[subsection]{Situation}

\theoremstyle{remark}
\newtheorem{remark}[subsection]{Remark}
\newtheorem{remarks}[subsection]{Remarks}

\numberwithin{equation}{subsection}

% Macros
%
\def\lim{\mathop{\mathrm{lim}}\nolimits}
\def\colim{\mathop{\mathrm{colim}}\nolimits}
\def\Spec{\mathop{\mathrm{Spec}}}
\def\Hom{\mathop{\mathrm{Hom}}\nolimits}
\def\Ext{\mathop{\mathrm{Ext}}\nolimits}
\def\SheafHom{\mathop{\mathcal{H}\!\mathit{om}}\nolimits}
\def\SheafExt{\mathop{\mathcal{E}\!\mathit{xt}}\nolimits}
\def\Sch{\mathit{Sch}}
\def\Mor{\mathop{\mathrm{Mor}}\nolimits}
\def\Ob{\mathop{\mathrm{Ob}}\nolimits}
\def\Sh{\mathop{\mathit{Sh}}\nolimits}
\def\NL{\mathop{N\!L}\nolimits}
\def\CH{\mathop{\mathrm{CH}}\nolimits}
\def\proetale{{pro\text{-}\acute{e}tale}}
\def\etale{{\acute{e}tale}}
\def\QCoh{\mathit{QCoh}}
\def\Ker{\mathop{\mathrm{Ker}}}
\def\Im{\mathop{\mathrm{Im}}}
\def\Coker{\mathop{\mathrm{Coker}}}
\def\Coim{\mathop{\mathrm{Coim}}}

% Boxtimes
%
\DeclareMathSymbol{\boxtimes}{\mathbin}{AMSa}{"02}

%
% Macros for moduli stacks/spaces
%
\def\QCohstack{\mathcal{QC}\!\mathit{oh}}
\def\Cohstack{\mathcal{C}\!\mathit{oh}}
\def\Spacesstack{\mathcal{S}\!\mathit{paces}}
\def\Quotfunctor{\mathrm{Quot}}
\def\Hilbfunctor{\mathrm{Hilb}}
\def\Curvesstack{\mathcal{C}\!\mathit{urves}}
\def\Polarizedstack{\mathcal{P}\!\mathit{olarized}}
\def\Complexesstack{\mathcal{C}\!\mathit{omplexes}}
% \Pic is the operator that assigns to X its picard group, usage \Pic(X)
% \Picardstack_{X/B} denotes the Picard stack of X over B
% \Picardfunctor_{X/B} denotes the Picard functor of X over B
\def\Pic{\mathop{\mathrm{Pic}}\nolimits}
\def\Picardstack{\mathcal{P}\!\mathit{ic}}
\def\Picardfunctor{\mathrm{Pic}}
\def\Deformationcategory{\mathcal{D}\!\mathit{ef}}
\usepackage{cleveref}
\crefname{originallemma}{original lemma}{original lemmas}
\Crefname{originallemma}{Original lemma}{Original lemmas}
\makeatletter
\newlabel{more-algebra-lemma-lift-projective-module}{{07M5}{source}{}{sourcepointer-1}{}}
\newlabel{more-algebra-lemma-lift-projective-module@cref}{{[originallemma][0][]07M5}{[1][0][]source}{}{}{}}
\makeatother
\crefname{originallemma}{original lemma}{original lemmas}
\Crefname{originallemma}{Original lemma}{Original lemmas}
\makeatletter
\newlabel{algebra-lemma-charpoly}{{00DX}{source}{}{sourcepointer-2}{}}
\newlabel{algebra-lemma-charpoly@cref}{{[originallemma][0][]00DX}{[1][0][]source}{}{}{}}
\makeatother
\crefname{originallemma}{original lemma}{original lemmas}
\Crefname{originallemma}{Original lemma}{Original lemmas}
\makeatletter
\newlabel{lemma-lift-idempotent-upstairs}{{07M4}{source}{}{sourcepointer-3}{}}
\newlabel{lemma-lift-idempotent-upstairs@cref}{{[originallemma][0][]07M4}{[1][0][]source}{}{}{}}
\makeatother
\title{A finite projective module modulo an ideal lifts after an etale ring extension}
\author{The Stacks Project Authors}
\date{stacks-project:more-algebra.tex:local-snapshot-20261012; GNU FDL 1.2 or later}
\begin{document}
\maketitle
\setcounter{subsection}{1}
\noindent\textit{Editorial note.} All delivered mathematical text is the author's own, copied from the pinned Stacks Project chapter file more-algebra.tex; no statement or proof step was rewritten, reordered or removed.\par\medskip
\noindent\textit{Editorial note.} The chapter file pinned-original-source.tex carries only \textbackslash{}input\{preamble\} on its first line, so no preamble text lives inside the chapter. The retained file sources/stacks-more-algebra/retained-preamble.tex is a mechanical, line-for-line copy of the project's single shared preamble as staged in the snapshot (Xy-pic setup plus theorem declarations and mathematical macros). Nothing was added, removed or redefined, and the excerpt serves only as preamble, never as a span coordinate.\par\medskip
\noindent\textit{Editorial note.} Both counted spans start at the \textbackslash{}begin line of their environment, so the statement and the author's proof are delivered whole; surrounding prose of the section is deliberately outside the counted spans and is carried as standard prerequisites instead.\par\medskip
\noindent\textit{Editorial note.} The citation fields (source\_id, source\_version, authors, source\_url, license.name, source\_title) are identical to those of every other card this lane has produced from this chapter, so the batch selector treats them as one source.\par\medskip
\noindent\textit{Allowed standard prerequisites (editorial).}\begin{itemize}
\item Cayley-Hamilton and the characteristic polynomial of an endomorphism of a finite free module. the author cites Algebra, Lemma algebra-lemma-charpoly at proof lines 2033-2035. Supplies the polynomial f whose roots control the idempotent.
\item Lifting of idempotents along an etale ring map that is an isomorphism modulo the ideal. the statement counted by the sibling card S1-algebra-lift-idempotent at 1784-1803, cited at proof lines 2058-2060. Produces the idempotent of the finite extension ring that is used to define the lifted module.
\item A finite projective module is a direct summand of a finite free module, so it is the image of an idempotent projector. used at proof lines 2028-2032. Supplies the projector whose lift produces the lifted module.
\end{itemize}\medskip
\ifdefined\hypertarget\hypertarget{sourcepointer-1}{}\fi
\noindent\textit{Source reference.} Original lemma 07M5: focused-seeds/AGP/sources/stacks-snapshot-20261011/tags.txt, source lines 9715--9715. The lemma's own canonical tag, quoted so that readers can resolve the label inside the delivered text; it points at the counted statement itself and carries no proof content.\par\medskip
\ifdefined\hypertarget\hypertarget{sourcepointer-2}{}\fi
\noindent\textit{Source reference.} Original lemma 00DX: focused-seeds/AGP/sources/stacks-snapshot-20261011/tags.txt, source lines 514--514. Navigation pointer to a statement the counted proof refers to but does not contain; nothing load-bearing is delegated to it.\par\medskip
\ifdefined\hypertarget\hypertarget{sourcepointer-3}{}\fi
\noindent\textit{Source reference.} Original lemma 07M4: focused-seeds/AGP/sources/stacks-snapshot-20261011/tags.txt, source lines 9714--9714. Navigation pointer to a statement the counted proof refers to but does not contain; nothing load-bearing is delegated to it.\par\medskip
\begin{lemma}
\label{lemma-lift-projective-module}
Let $A$ be a ring, let $I \subset A$ be an ideal.
Let $\overline{P}$ be a finite projective $A/I$-module.
Then there exists an \'etale ring map $A \to A'$ which induces
an isomorphism $A/I \to A'/IA'$ and a finite projective
$A'$-module $P'$ lifting $\overline{P}$.
\end{lemma}

\begin{proof}
We can choose an integer $n$ and a direct sum decomposition
$(A/I)^{\oplus n} = \overline{P} \oplus \overline{K}$
for some $R/I$-module $\overline{K}$. Choose a lift
$\varphi : A^{\oplus n} \to A^{\oplus n}$ of the projector $\overline{p}$
associated to the direct summand $\overline{P}$.
Let $f \in A[x]$ be the characteristic polynomial of $\varphi$.
Set $B = A[x]/(f)$. By Cayley-Hamilton
(Algebra, Lemma \ref{algebra-lemma-charpoly}) there is a map
$B \to \text{End}_A(A^{\oplus n})$ mapping $x$ to $\varphi$.
For every prime $\mathfrak p \supset I$ the image of $f$ in
$\kappa(\mathfrak p)$ is $(x - 1)^rx^{n - r}$ where $r$ is the
dimension of $\overline{P} \otimes_{A/I} \kappa(\mathfrak p)$.
Hence $(x - 1)^nx^n$ maps to zero in $B \otimes_A \kappa(\mathfrak p)$
for all $\mathfrak p \supset I$. Thus $x(1 - x)$ is contained
in every prime ideal of $B/IB$. Hence $x^N(1 - x)^N$ is
contained in $IB$ for some $N \geq 1$.
It follows that $x^N + (1 - x)^N$ is a unit in $B/IB$ and that
$$
\overline{e} = \text{image of }\frac{x^N}{x^N + (1 - x)^N}\text{ in }B/IB
$$
is an idempotent as both assertions hold in $\mathbf{Z}[x]/(x^N(x - 1)^N)$.
The image of $\overline{e}$ in $\text{End}_{A/I}((A/I)^{\oplus n})$ is
$$
\frac{\overline{p}^N}{\overline{p}^N + (1 - \overline{p})^N} = \overline{p}
$$
as $\overline{p}$ is an idempotent. After replacing $A$ by an \'etale
extension $A'$ as in the lemma, we may assume there exists an idempotent
$e \in B$ which maps to $\overline{e}$ in $B/IB$, see
Lemma \ref{lemma-lift-idempotent-upstairs}.
Then the image of $e$ under the map
$$
B = A[x]/(f) \longrightarrow \text{End}_A(A^{\oplus n}).
$$
is an idempotent element $p$ which lifts $\overline{p}$.
Setting $P = \Im(p)$ we win.
\end{proof}

\end{document}
